\section*{Bab \ref{ch:g10:refraction} --- Pembiasan Cahaya}

\begin{solution}{exo:g10:refraction:1}
Sudut diukur dari \omterm{def:g10:refraction:angles}{garis normal}: \omterm{def:g10:refraction:angles}{sudut datangnya}
$90^\circ - 25^\circ = 65^\circ$, sehingga sudut pantulnya juga
$65^\circ$. Berkas datang dan berkas pantulnya membentuk sudut
$2 \times 65^\circ = 130^\circ$ satu sama lain.
\end{solution}

\begin{solution}{exo:g10:refraction:2}
\emph{1.} $v = c/n$: di dalam air
$v = \qty{3.00e8}{m/s}/1.33 = \qty{2.26e8}{m/s}$; di dalam intan
$v = \qty{3.00e8}{m/s}/2.42 = \qty{1.24e8}{m/s}$.

\emph{2.} $n = \dfrac{c}{v} = \dfrac{\qty{3.00e8}{m/s}}
{\qty{1.97e8}{m/s}} = 1.52$: kaca.
\end{solution}

\begin{solution}{exo:g10:refraction:3}
$\sin i_2 = \dfrac{\sin 40^\circ}{1.50} = \dfrac{0.643}{1.50} = 0.429$,
sehingga $i_2 = 25.4^\circ$. Karena lintasan cahaya dapat dibalik,
sinar di dalam kaca yang mengenai permukaan pada $25.4^\circ$ keluar ke
udara tepat pada $40^\circ$.
\end{solution}

\begin{solution}{exo:g10:refraction:4}
$\sin i_2 = \dfrac{\sin 55^\circ}{1.33} = \dfrac{0.819}{1.33} = 0.616$,
sehingga $i_2 = 38.0^\circ$. Sinarnya membelok \emph{mendekati} \omterm{def:g10:refraction:angles}{garis
normal}, karena air adalah medium yang lebih lambat ($n_2 > n_1$).
\end{solution}

\begin{solution}{exo:g10:refraction:5}
$n = \dfrac{\sin 45^\circ}{\sin 32^\circ} = \dfrac{0.707}{0.530}
= 1.33$: zat cair itu adalah air.
\end{solution}

\begin{solution}{exo:g10:refraction:6}
\emph{1.} Kelima nisbahnya adalah
\[
\frac{0.342}{0.233} = 1.47,\quad
\frac{0.500}{0.334} = 1.50,\quad
\frac{0.643}{0.431} = 1.49,\quad
\frac{0.766}{0.515} = 1.49,\quad
\frac{0.866}{0.581} = 1.49 :
\]
tetap dalam batas galat pengukuran --- itulah \omterm{thm:g10:refraction:snell}{hukum Snell}.

\emph{2.} $n \approx 1.49$: pleksiglas.

\emph{3.} \omterm{thm:g10:refraction:snell}{Hukum Snell} meramalkan $\sin i_2 = \sin i_1 / n$: digambarkan
sinus terhadap sinus, datanya harus berjajar pada garis lurus melalui
titik asal, dan sebuah penggaris sudah dapat memeriksanya sekali
pandang. Digambarkan sudut terhadap sudut, datanya terletak pada sebuah
lengkung; pada sudut kecil lengkung itu sangat mirip garis lurus, dan
justru itulah yang mengelabui Ptolemaios.
\end{solution}

\begin{solution}{exo:g10:refraction:7}
$\sin i_c = n_2/n_1$ dengan $n_2 = 1.00$:
air $\sin i_c = 1/1.33 = 0.752$, $i_c = 48.8^\circ$;
kaca $\sin i_c = 1/1.50 = 0.667$, $i_c = 41.8^\circ$;
intan $\sin i_c = 1/2.42 = 0.413$, $i_c = 24.4^\circ$.

Dari udara ke dalam air, $\sin i_2 = \sin i_1/1.33 \leq 1/1.33 < 1$
untuk setiap \omterm{def:g10:refraction:angles}{sudut datang}: sudut biasnya selalu ada, sehingga tidak
pernah terjadi apa pun yang istimewa --- \omterm{def:g10:refraction:critical}{pemantulan sempurna} menuntut
$n_1 > n_2$.
\end{solution}

\begin{solution}{exo:g10:refraction:8}
\emph{1.} $d' = \dfrac{d}{n} = \dfrac{\qty{2.0}{m}}{1.33}
= \qty{1.5}{m}$.

\emph{2.} Ikan itu tampak pada $\qty{40}{cm}/1.33 = \qty{30}{cm}$:
bayangannya berada \emph{di atas} ikannya, sehingga bangau itu harus
mematuk \emph{di bawah} bayangan yang dilihatnya. Bangau, yang tidak
terbebani \omterm{thm:g10:refraction:snell}{hukum Snell}, mempelajarinya dengan mencoba-coba.
\end{solution}

\begin{solution}{exo:g10:refraction:9}
\emph{1.} $\sin r = \dfrac{\sin 50^\circ}{1.50} = \dfrac{0.766}{1.50}
= 0.511$, sehingga $r = 30.7^\circ$.

\emph{2.} Dari kaca ke udara pada $30.7^\circ$:
$\sin i_2 = 1.50 \times 0.511 = 0.766$, sehingga $i_2 = 50^\circ$ ---
sudut keluarnya sama dengan sudut masuknya.

\emph{3.} Kedua \omterm{def:g10:refraction:refraction}{pembiasan} itu saling meniadakan: sinarnya pergi
\emph{sejajar} dengan arah semula, hanya bergeser ke samping. Itulah
sebabnya kaca jendela tidak mengubah bentuk pemandangan, melainkan
hanya menggesernya (tanpa terasa).
\end{solution}

\begin{solution}{exo:g10:refraction:10}
$\sin 55^\circ = 0.819$. Merah:
$\sin i_2 = 0.819/1.61 = 0.509$, $i_2 = 30.6^\circ$. Biru:
$\sin i_2 = 0.819/1.66 = 0.493$, $i_2 = 29.6^\circ$. Kedua warna itu
terpisah sebesar $1.0^\circ$; biru, yang indeksnya lebih besar,
membelok lebih kuat dan berakhir lebih dekat dengan \omterm{def:g10:refraction:angles}{garis normal}.
\end{solution}

\begin{solution}{exo:g10:refraction:11}
\emph{1.} $\sin i_c = \dfrac{1.46}{1.48} = 0.986$, sehingga
$i_c = 80.6^\circ$.

\emph{2.} Sinar pada $i_c$ dari \omterm{def:g10:refraction:angles}{garis normal} membentuk $90^\circ - i_c$
dengan sumbunya, sehingga lintasannya sepanjang sumbu $L$ adalah
$\dfrac{L}{\sin i_c}$:
\[
\qty{1000}{m} \times \frac{1.48}{1.46} = \qty{1013.7}{m},
\]
yaitu sekitar \qty{14}{m} lebih panjang daripada sinar yang menyusuri
sumbunya.
\end{solution}

\begin{solution}{exo:g10:refraction:12}
\emph{1.} Cahaya yang menyusur datang dari kaki langit tiba pada $i_1$
yang mendekati $90^\circ$; cahaya itu dibiaskan menjadi
$\sin i_2 = \sin 90^\circ / 1.33 = 0.752$, yaitu $i_2 = 48.8^\circ$ ---
sudut kritisnya, sebagaimana dituntut sifat dapat dibaliknya lintasan
cahaya. Karena itu sinar dari seluruh langit mencapai mata di dalam
kerucut bersetengah sudut $48.8^\circ$.

\emph{2.} Tepi cakramnya berada di tempat kerucut itu menemui
permukaan:
\[
r = d \tan i_c = \qty{2.0}{m} \times \tan 48.8^\circ
= \qty{2.0}{m} \times 1.14 = \qty{2.3}{m}.
\]

\emph{3.} Di luar cakram itu, cahaya dari bawah mengenai permukaan di
luar \omterm{def:g10:refraction:critical}{sudut kritis} dan dipantulkan sempurna: penyelam itu melihat cermin
keperakan yang menampilkan dasar kolamnya.
\end{solution}

\begin{solution}{exo:g10:refraction:13}
\emph{1.} Merah: $\sin i_2 = 1.51 \times \sin 30^\circ = 0.755$,
sehingga $i_2 = 49.0^\circ$. Biru: $\sin i_2 = 1.53 \times 0.500 = 0.765$, sehingga $i_2 = 49.9^\circ$. Berkas putih itu terkipas
sebesar $49.9^\circ - 49.0^\circ = 0.9^\circ$ --- \omterm{def:g10:light-spectra:white}{spektrum} prisma dari
satu kali \omterm{def:g10:refraction:refraction}{pembiasan}.

\emph{2.} Pada $45^\circ$: $\sin i_2 = 1.51 \times 0.707 = 1.07 > 1$
(dan $1.53 \times 0.707 = 1.08$ untuk biru). Tidak ada sudut bias bagi
kedua warna itu: sudut kritisnya
$\arcsin(1/1.51) = 41.5^\circ < 45^\circ$, sehingga sisi panjangnya
memantulkan seluruh berkas dengan sempurna.
\end{solution}

\begin{solution}{exo:g10:refraction:14}
\emph{1.} \omterm{def:g10:refraction:critical}{Sudut kritis} kaca--udara: $\sin i_c = 1/1.50$, sehingga
$i_c = 41.8^\circ < 45^\circ$: berkasnya dipantulkan sempurna ---
$100\%$ cahayanya, tanpa lapisan yang dapat memudar. Cermin logam
kehilangan beberapa persen pada setiap pantulan, dan sebuah periskop
memerlukan dua pantulan.

\emph{2.} Kaca--air: $\sin i_c = \dfrac{1.33}{1.50} = 0.887$, sehingga
$i_c = 62.5^\circ$. Kini $45^\circ < i_c$: \omterm{def:g10:refraction:angles}{sudut datangnya} di bawah
\omterm{def:g10:refraction:critical}{sudut kritis}, sebagian besar cahayanya dibiaskan keluar ke air, dan
cerminnya --- karena itu juga bayangannya --- gagal.

\emph{3.} $\sin i_2 = \dfrac{1.50}{1.33} \times \sin 45^\circ
= 1.128 \times 0.707 = 0.798$, sehingga $i_2 = 52.9^\circ$ ke dalam air.
\end{solution}

\begin{solution}{exo:g10:refraction:15}
\emph{1.} $v = \dfrac{\qty{3.00e8}{m/s}}{1.47} = \qty{2.04e8}{m/s}$,
sehingga
\[
t = \frac{\qty{6.2e6}{m}}{\qty{2.04e8}{m/s}} = \qty{30.4}{ms} \approx
\qty{30}{ms}.
\]

\emph{2.} Radio: $t = \qty{6.2e6}{m}/\qty{3.00e8}{m/s} =
\qty{20.7}{ms}$. Kaca itu berharga $30.4 - 20.7 = \qty{9.7}{ms}$ pada
tiap arah, yaitu sekitar \qty{19}{ms} untuk sekali pulang pergi.

\emph{3.} Turunkan indeksnya --- serat berinti rongga memandu cahaya
melalui udara ($n \approx 1.00$), dan sambungan gelombang mikro melalui
atmosfer melakukan hal yang sama --- atau perpendek lintasannya: kabel
yang dipasang lebih dekat dengan rute lingkaran besar antara kedua kota
itu. Keduanya benar-benar diperjualbelikan, dengan harga per milisekon.
\end{solution}

\begin{solution}{pb:g10:refraction:1}
\textbf{1.} $v_{\text{air}} = \qty{3.00e8}{m/s}/1.33 =
\qty{2.26e8}{m/s}$; $v_{\text{intan}} = \qty{3.00e8}{m/s}/2.42 =
\qty{1.24e8}{m/s}$ --- intan memperlambat cahaya dengan faktor
$n = 2.42$.

\textbf{2.} $\sin i_2 = \sin 40^\circ / 1.33 = 0.643/1.33 = 0.483$,
sehingga $i_2 = 28.9^\circ$.

\textbf{3.} $\sin i_2 = 1.33 \times \sin 28.9^\circ = 1.33 \times
0.483 = 0.643$, sehingga $i_2 = 40^\circ$: cahaya menelusuri kembali
lintasannya --- sifat dapat dibaliknya sinar cahaya.

\textbf{4.} $\sin i_c = 1/1.33 = 0.752$, sehingga $i_c = 48.8^\circ$.
Pada $60^\circ > i_c$, \omterm{thm:g10:refraction:snell}{hukum Snell} akan menuntut $\sin i_2 = 1.33 \sin
60^\circ = 1.15 > 1$: tidak ada \omterm{def:g10:refraction:refraction}{sinar bias}; permukaannya memantulkan
seluruh cahaya kembali ke bawah.

\textbf{5.} $v = \qty{3.00e8}{m/s}/1.0003 = \qty{2.999e8}{m/s}$: cahaya
di udara lebih lambat daripada di ruang hampa sebesar $0.03\%$, jauh di
bawah ketelitian pengukuran sudut kita, sehingga
$n_{\text{udara}} = 1.00$ tidak merugikan.

\textbf{6.} $d' = \qty{3.0}{m}/1.33 = \qty{2.3}{m}$: kolam itu
menyembunyikan \qty{70}{cm} dari dirinya sendiri.

\textbf{7.} Dua sinar yang meninggalkan uang logam itu dibiaskan pada
permukaan, membelok \emph{menjauhi} \omterm{def:g10:refraction:angles}{garis normal}, sehingga keduanya
memencar lebih tajam di atas air daripada di bawahnya. Mata
memperpanjang kedua sinar yang keluar itu ke belakang sebagai garis
lurus; keduanya berpotongan \emph{di atas} uang logam itu, dan di
titik potong itulah mata menempatkan bayangannya.

\textbf{8.} $r = d \tan i_c = \qty{3.0}{m} \times \tan 48.8^\circ =
\qty{3.0}{m} \times 1.14 = \qty{3.4}{m}$: cakram selebar sekitar
\qty{6.9}{m} memuat seluruh langit.

\textbf{9.} Pada $30^\circ$: $\sin i_2 = 1.33 \times 0.500 = 0.665$,
sinarnya keluar pada $41.7^\circ$. Pada $45^\circ$: $\sin i_2 = 1.33 \times
0.707 = 0.940$, sinarnya keluar pada $70.1^\circ$, hampir
menyusur. Pada $50^\circ$: $1.33 \times 0.766 = 1.02 > 1$ ---
\omterm{def:g10:refraction:critical}{pemantulan sempurna}, sinarnya menyelam kembali ke dalam kolam.

\textbf{10.} Bayangannya berada di atas ikan yang sebenarnya
(\cref{prop:g10:refraction:depth}), jadi bidiklah \emph{di bawahnya}.
Dan benar: ikan itu juga melihat si penombak tergeser, termampatkan ke
arah tegak di dalam jendela Snell --- masing-masing melihat yang lain
di tempat yang bukan tempatnya.

\textbf{11.} Intan: $\sin i_c = 1/2.42 = 0.413$, sehingga
$i_c = 24.4^\circ$, berbanding $41.1^\circ$ untuk kaca: kerucut keluar
intan bahkan tidak selebar separuhnya.

\textbf{12.} Di dalam intan, $30^\circ > 24.4^\circ$: \omterm{def:g10:refraction:critical}{pemantulan
sempurna}, sinarnya tetap tinggal di dalam batu. Di dalam kaca,
$30^\circ < 41.1^\circ$: sinarnya bocor keluar lewat belakang. Faset
potongan brilian ditata bersudut sedemikian rupa sehingga cahaya yang
masuk hampir selalu mengenainya di luar $24.4^\circ$ dan memantul-mantul
sampai keluar ke atas menuju mata --- tiruan dari kaca, dengan kerucut
keluarnya yang lebar, begitu saja meloloskan cahayanya.

\textbf{13.} $\sin 45^\circ = 0.707$. Merah: $0.707/2.41 = 0.293$,
sehingga $i_2 = 17.1^\circ$. Biru: $0.707/2.45 = 0.289$, sehingga
$i_2 = 16.8^\circ$. Pemisahan pada saat masuk: sekitar $0.3^\circ$.

\textbf{14.} $\sin 17.0^\circ = 0.292$. Merah: $\sin i_2 = 2.41 \times
0.292 = 0.705$, sehingga $i_2 = 44.8^\circ$. Biru: $\sin i_2 = 2.45 \times
0.292 = 0.716$, sehingga $i_2 = 45.8^\circ$. Kipasnya kini
selebar $1.0^\circ$: \omterm{def:g10:refraction:refraction}{pembiasan} keluarnya, yang bekerja di dekat \omterm{def:g10:refraction:critical}{sudut
kritis}, meregangkan pemisahan masuk $0.3^\circ$ itu menjadi tiga kali
lipat. Semburan warna yang teregang itulah \emph{api} sang intan.

\textbf{15.} Intan--air: $\sin i_c = 1.33/2.42 = 0.550$, sehingga
$i_c = 33.3^\circ$. Sinar $30^\circ$ yang tadinya terperangkap ketika
fasetnya menghadap udara kini lolos ($30^\circ < 33.3^\circ$) ke dalam
lapisan airnya: setiap tetes, sidik jari atau lapisan minyak menaikkan
indeks di luarnya, melebarkan kerucut keluarnya dan menguras kilaunya
--- itulah gunanya kain para pengasah permata.

\textbf{16.} $\sin i_c = \dfrac{1.46}{1.48} = 0.986$, sehingga
$i_c = 80.6^\circ$.

\textbf{17.} \omterm{def:g10:refraction:critical}{Pemantulan sempurna} terjadi \emph{pada permukaannya},
sehingga permukaan itu harus tetap sempurna: pada benang telanjang,
setiap butir debu, goresan, tetes air atau jari yang menyentuh akan
menaikkan indeks di luar titik itu dan membocorkan cahayanya. \omterm{def:g10:refraction:fiber}{Selubung}
memendam cerminnya di dalam kaca yang bersih, tempat tidak ada apa pun
yang dapat mencapainya --- dan, sebagai imbalan tambahan, melindungi
\omterm{def:g10:refraction:fiber}{inti} setipis rambut itu secara mekanis.

\textbf{18.} $v = \qty{3.00e8}{m/s}/1.48 = \qty{2.03e8}{m/s}$, sehingga
\[
t = \frac{\qty{6.0e6}{m}}{\qty{2.03e8}{m/s}} = \qty{29.6}{ms} \approx
\qty{30}{ms}.
\]

\textbf{19.} Faktor lintasannya $1/\sin i_c = 1.48/1.46 = 1.0137$: pada
kasus terburuk $\qty{6000}{km} \times 0.0137 \approx \qty{82}{km}$ kaca
tambahan, yang berharga $\qty{8.2e4}{m} / \qty{2.03e8}{m/s} \approx \qty{0.4}{ms}$ --- zigzagnya nyaris cuma-cuma.

\textbf{20.} Ping: $2 \times (29.6 + 0.4) \approx \qty{60}{ms}$.
Kapasitasnya: $\dfrac{10^{13}}{2.4 \times 10^8} \approx 40\,000$
salinan buku ini tiap sekon. Dalam satu kalimat: dengan memerangkap
cahaya di dalam kaca melalui \omterm{def:g10:refraction:critical}{pemantulan sempurna}, \omterm{thm:g10:refraction:snell}{hukum Snell}
menyeberangkan satu perpustakaan melintasi samudra setiap sekon, enam
puluh milisekon pergi dan kembali.
\end{solution}
